\documentclass[11pt]{amsart}

\usepackage[margin=1in]{geometry}
\usepackage{amsmath,amssymb,amsthm}
\usepackage{mathtools}
\usepackage[colorlinks=true,linkcolor=blue,citecolor=blue,urlcolor=blue]{hyperref}
\usepackage{enumitem}

% --- theorem environments ---
\theoremstyle{plain}
\newtheorem{theorem}{Theorem}[section]
\newtheorem{proposition}[theorem]{Proposition}
\newtheorem{lemma}[theorem]{Lemma}
\newtheorem{corollary}[theorem]{Corollary}

\theoremstyle{definition}
\newtheorem{definition}[theorem]{Definition}
\newtheorem{example}[theorem]{Example}
\newtheorem{fact}[theorem]{Fact}
\newtheorem{problem}[theorem]{Problem}

\theoremstyle{remark}
\newtheorem{remark}[theorem]{Remark}

% --- macros ---
\newcommand{\R}{\mathbb{R}}
\newcommand{\dom}{\operatorname{dom}}
\newcommand{\intr}{\operatorname{int}}
\newcommand{\cl}{\operatorname{cl}}
\newcommand{\D}{D_f}
\newcommand{\Proj}{\overrightarrow{\Pi}}
\newcommand{\Cstar}{C^{*}}
\newcommand{\ip}[2]{\langle #1, #2 \rangle}

\title[Full domain in a theorem on right Bregman--Chebyshev sets]
{Full domain cannot be dropped from a theorem\\ on right Bregman--Chebyshev sets}

\author{}

\thanks{This investigation was carried out with substantial assistance from
automated systems, and we omit a conventional byline rather than overstate the
named researchers' personal contributions.
\emph{Nikol Savova}, University of Oxford,
\href{mailto:nikol.p.savova@gmail.com}{\texttt{nikol.p.savova@gmail.com}};
\emph{Sihao Huang}, independent researcher,
\href{mailto:sihao.c.huang@gmail.com}{\texttt{sihao.c.huang@gmail.com}}.}

\date{\today}

\begin{document}

\begin{abstract}
Bauschke, Macklem and Wang record the following fact: if $f$ is a Legendre
function on $X$ with $\dom f = X$, if $C \subseteq U = \intr \dom f$ is closed,
nonempty and satisfies $\cl \nabla f(C) \subseteq U^{*}$, and if $C$ is right
$\D$-Chebyshev, then $\nabla f(C)$ is convex. They ask whether the hypothesis
$\dom f = X$ is necessary. We show that it cannot be deleted. Taking $f$ to be
the negative entropy on $\R^{2}$ and
$C = \{(e^{t}, e^{-t^{2}}) : t \in [1,2]\}$, the set $C$ is compact and right
$\D$-Chebyshev at every point of $U$, the closure hypothesis
$\cl \nabla f(C) \subseteq U^{*}$ holds, and $\nabla f(C)$ is a strictly
concave arc, hence nonconvex. The verification is elementary: the right
Bregman projection onto $C$ reduces to minimising a one-variable function whose
second derivative is a sum of manifestly positive terms on $[1,2]$.

The same example bears on later work of Luo, Meng, Wen and Yao. Their
Theorem~3.12(2) asserts, for $U = X$, the equivalence of a variational condition
(i) with the property (ii) that $C$ is a right $\D$-sun. We show by direct
computation that our $C$ is a right $\D$-sun while (i) fails for an explicit
point, so the requirement $U = X$ cannot be deleted from their theorem either.
The obstruction is visible: the ray along which the sun property is applied
leaves $U$ after a bounded distance.

We also record, with attribution, that the remaining hypothesis
$\cl \nabla f(C) \subseteq U^{*}$ is redundant, as a short consequence of the
results of Luo, Meng, Wen and Yao.
\end{abstract}

\maketitle

\section{Introduction}

Let $X = \R^{n}$ with the standard inner product, and let
$f \colon X \to \left]-\infty, +\infty\right]$ be proper, lower semicontinuous
and convex. Write $U = \intr \dom f$ and $U^{*} = \intr \dom f^{*}$, where
$f^{*}$ is the Fenchel conjugate. Recall that $f$ is \emph{Legendre} if it is
both essentially smooth and essentially strictly convex \cite{BBC2001,Rock1970}.
For $x \in \dom f$ and $y \in U$, the \emph{Bregman distance} is
\[
  \D(x,y) = f(x) - f(y) - \ip{\nabla f(y)}{x - y}.
\]
It is not symmetric, so a set has two projection operators. Throughout, we use
the \emph{right} one, in which the second argument varies:
\[
  \Proj_{C}(x) = \operatorname*{arg\,min}_{y \in C} \D(x,y),
  \qquad x \in U .
\]
Following \cite{BMW2010} we call $C$ \emph{right $\D$-Chebyshev} if
$\Proj_{C}(x)$ is a singleton for every $x \in U$. We abbreviate
\[
  \Cstar := \nabla f(C).
\]

The starting point of this note is the following statement, recorded as
Fact~3.2 of the survey of Bauschke, Macklem and Wang \cite{BMW2010}.

\begin{fact}[{\cite[Fact 3.2]{BMW2010}}]\label{fact:BMW}
Suppose that
\begin{enumerate}[label=\textup{(\alph*)},leftmargin=2.2em]
  \item $\dom f = X$;
  \item $C \subseteq U$ is closed and nonempty, and $\cl \Cstar \subseteq U^{*}$;
  \item $C$ is right $\D$-Chebyshev.
\end{enumerate}
Then $\Cstar$ is convex.
\end{fact}

The survey attaches an open problem to this fact, which we restate.

\begin{problem}[{\cite[Problem 2]{BMW2010}}]\label{prob:2}
Is the hypothesis $\dom f = X$ necessary? Either exhibit a Legendre $f$ with
$\dom f \neq X$ and a closed nonempty right $\D$-Chebyshev $C \subseteq U$ with
$\Cstar$ nonconvex, or show that the conclusion survives without full domain.
\end{problem}

We settle Problem~\ref{prob:2} in the first of the two directions.

\begin{theorem}\label{thm:main}
Let $X = \R^{2}$ and let $f$ be the negative entropy,
\[
  f(x) = g(x_{1}) + g(x_{2}),
  \qquad
  g(s) =
  \begin{cases}
    s \ln s - s, & s > 0,\\
    0, & s = 0,\\
    +\infty, & s < 0.
  \end{cases}
\]
Fix real numbers $a, b$ with $1/\sqrt{2} \le a < b$ and put
\[
  C = \bigl\{ (e^{t}, e^{-t^{2}}) : t \in [a,b] \bigr\} .
\]
Then $f$ is Legendre with $\dom f = \R^{2}_{+} \neq X$, and:
\begin{enumerate}[label=\textup{(\roman*)},leftmargin=2.2em]
  \item $C$ is a nonempty compact subset of $U = \R^{2}_{++}$;
  \item $\cl \Cstar = \Cstar \subseteq U^{*} = \R^{2}$, so hypothesis \textup{(b)} of
        Fact~\textup{\ref{fact:BMW}} holds;
  \item $C$ is right $\D$-Chebyshev, so hypothesis \textup{(c)} holds;
  \item $\Cstar = \{ (t, -t^{2}) : t \in [a,b] \}$ is not convex.
\end{enumerate}
Consequently the implication of Fact~\textup{\ref{fact:BMW}} becomes false if
hypothesis \textup{(a)} is deleted while \textup{(b)} and \textup{(c)} are retained.
\end{theorem}

Two features of the example are worth stating at the outset, because both are
easy to lose in a more elaborate construction.

First, hypothesis (b) is genuinely satisfied, not merely arranged to be
vacuous. If $\cl \Cstar \subseteq U^{*}$ failed, the example would be deleting
two hypotheses at once and would say nothing about (a) in particular. Here
$U^{*} = \R^{2}$ and $\Cstar$ is compact, so (b) holds for the strongest reason
available.

Second, $C$ is right $\D$-Chebyshev at \emph{every} point of $U$, including
points arbitrarily close to $\partial U$ and points far from $C$. Uniqueness is
not a local or generic phenomenon here. It follows from a single inequality
(Lemma~\ref{lem:convex}) which holds uniformly in $x \in U$.

The verification is elementary throughout, and the proofs below use no numerical
computation. The one inequality that carries the argument is
\[
  4t^{2} - 2 \ \ge\ 0
  \qquad\text{for } t \ge 1/\sqrt{2},
\]
after which the required second derivative is a sum of nonnegative terms plus a
strictly positive one. We draw attention to this because it is the reason the
example is checkable by hand, and because it explains the constraint
$a \ge 1/\sqrt{2}$, which is the only role played by the choice of interval.

Section~\ref{sec:luo} relates the example to subsequent work of Luo, Meng, Wen
and Yao \cite{LMWY2019}, where it turns out to have a second use, and
Section~\ref{sec:b} records what is known about the remaining hypothesis (b).
Section~\ref{sec:open} states what we do \emph{not} claim, and what remains open.

\subsection*{Relation to earlier examples, and what is not claimed}

The example is a variation on a template that is not ours. Bauschke, Wang, Ye
and Yuan \cite[Example 7.5]{BWYY2009}, reproduced as Example~3.3 of the survey
\cite{BMW2010}, work on the exponential side with
$h(u) = e^{u_{1}} + e^{u_{2}}$ and the segment
$A = \{ (\lambda, 2\lambda) : 0 \le \lambda \le 1 \}$, whose image under
$\nabla h$ is the nonconvex curve $\{ (e^{\lambda}, e^{2\lambda}) \}$. Since the
right projection for $f$ corresponds to the left projection for $f^{*}$
\cite[Lemma 3.2]{LMWY2019}, that curve is a nonconvex right $\D$-Chebyshev set
for the negative entropy, and full domain already fails there. But its gradient
image is $A$ itself, which is \emph{convex}: the conclusion of
Fact~\ref{fact:BMW} survives, and full domain is not challenged. That is exactly
why Problem~\ref{prob:2} was worth asking.

Our construction changes one thing. Where \cite[Example 7.5]{BWYY2009} takes the
dual set to be a convex segment, we take it to be a nonconvex parabolic arc, and
the work is in checking that global Chebyshevness survives the change.

\begin{center}
\begin{tabular}{lll}
\hline
 & \cite[Example 7.5]{BWYY2009} & Theorem~\ref{thm:main} \\
\hline
dual set $A$ & $\{ (\lambda, 2\lambda) \}$ & $\{ (t, -t^{2}) \}$ \\
$A$ convex? & yes & no \\
$C = \exp(A)$ right $\D$-Chebyshev? & yes & yes \\
$\Cstar = \nabla f(C) = A$ convex? & yes & no \\
bears on Problem~\ref{prob:2}? & no & yes \\
\hline
\end{tabular}
\end{center}

Stated on the dual side, Theorem~\ref{thm:main} says that the parabolic arc
$\{ (t,-t^{2}) : t \in [a,b] \}$ is a compact nonconvex set that is globally
left Chebyshev for the Bregman distance of $h(u) = e^{u_{1}} + e^{u_{2}}$. This
is the more transparent phrasing, and it makes clear that the obstruction lives
entirely on the dual side. It is worth noting that $h$ is not supercoercive.
Since $\dom f = X$ is equivalent to supercoercivity of $f^{*}$ for closed convex
$f$, the dual reading of Theorem~\ref{thm:main} is that a supercoercivity
hypothesis cannot be deleted from the corresponding left-projection statement
either.

We therefore claim neither the first nonconvex right $\D$-Chebyshev set for the
negative entropy --- that is \cite[Example 7.5]{BWYY2009} --- nor the duality
template, nor any new theory. Laude, Ochs and Cremers \cite{LOC2020} also give
nonconvex examples for entropy-like kernels, though their single-valuedness is
local, holding on a neighbourhood under relative prox-regularity rather than at
every point of $U$; they record Problem~\ref{prob:2} as open. Themelis and Wang
\cite{TW2025} develop Bregman operators on their natural domains, which is the
setting in which a statement like Theorem~\ref{thm:main} is naturally expressed.

The point of Theorem~\ref{thm:main} is narrower than any of these, and is
exactly what Problem~\ref{prob:2} asks for: a set that is right $\D$-Chebyshev
\emph{globally on $U$}, with \emph{nonconvex} gradient image, \emph{while
hypothesis \textup{(b)} holds}, so that full domain is the only hypothesis of
Fact~\ref{fact:BMW} that fails. We are not aware of an earlier example with this
combination of properties.

We also emphasise what Theorem~\ref{thm:main} does not say. It does not say
that $\dom f = X$ is necessary for $\Cstar$ to be convex in any individual
instance; convexity of $\Cstar$ certainly occurs for many $f$ of non-full
domain, \cite[Example 7.5]{BWYY2009} being one. It says that $\dom f = X$ cannot
be deleted from the \emph{universal} statement of Fact~\ref{fact:BMW}. Nor does
it identify the weakest hypothesis that could replace full domain; see
Section~\ref{sec:open}.

\section{The example}\label{sec:ex}

Throughout this section $f$, $C$, $a$ and $b$ are as in
Theorem~\ref{thm:main}, and we write
\[
  c(t) = (e^{t}, e^{-t^{2}}), \qquad t \in [a,b],
\]
so that $C = c([a,b])$. All closures are taken in $X = \R^{2}$.

\begin{lemma}\label{lem:legendre}
The function $f$ is Legendre, with
\[
  \dom f = \R^{2}_{+},
  \quad
  U = \R^{2}_{++},
  \quad
  \nabla f(x) = (\ln x_{1}, \ln x_{2}) \ \text{ on } U,
\]
and
\[
  f^{*}(u) = e^{u_{1}} + e^{u_{2}},
  \quad
  \dom f^{*} = U^{*} = \R^{2}.
\]
Moreover $\nabla f \colon U \to U^{*}$ is a bijection with inverse
$u \mapsto (e^{u_{1}}, e^{u_{2}})$.
\end{lemma}

\begin{proof}
For $s > 0$ we have $g'(s) = \ln s$ and $g''(s) = 1/s > 0$, so $g$ is strictly
convex on $\left]0,\infty\right[$. Strict convexity persists at the left
endpoint: for $b_{0} > 0$ and $0 < \lambda < 1$,
\[
  g(\lambda b_{0})
  = \lambda g(b_{0}) + \lambda b_{0} \ln \lambda
  < \lambda g(b_{0}) + (1 - \lambda) g(0),
\]
since $\ln \lambda < 0$ and $g(0) = 0$. As
$\lim_{s \downarrow 0} (s \ln s - s) = 0 = g(0)$, the function $g$ is lower
semicontinuous on $\R$. Hence $f$ is proper, lower semicontinuous and convex,
with $\dom f = \R^{2}_{+}$ and $U = \R^{2}_{++}$. Since
$(-1,0) \notin \dom f$, we have $\dom f \neq X$.

On $U$ the gradient is $\nabla f(x) = (\ln x_{1}, \ln x_{2})$, and the Hessian
$\operatorname{diag}(1/x_{1}, 1/x_{2})$ is positive definite. If
$x^{k} \in U$ converges to a boundary point of $\dom f$ then some coordinate
tends to $0$, so the corresponding component of $\nabla f(x^{k})$ tends to
$-\infty$ and $\lVert \nabla f(x^{k}) \rVert \to \infty$; thus $f$ is
essentially smooth.

No boundary point of $\dom f$ carries a subgradient. Indeed, suppose
$x_{j} = 0$ and $p \in \partial f(x)$. Perturbing only the $j$-th coordinate to
$s > 0$ in the subgradient inequality gives $s \ln s - s \ge p_{j} s$, hence
$\ln s - 1 \ge p_{j}$ for every $s > 0$; taking $s = e^{p_{j}}$ yields
$p_{j} - 1 \ge p_{j}$, a contradiction. Therefore
$\dom \partial f = U$, and since $f$ is strictly convex on $U$ it is
essentially strictly convex. So $f$ is Legendre.

For the conjugate, fix $u \in \R$ and consider $s \mapsto us - g(s)$ on
$[0,\infty[$. Its derivative on $\left]0,\infty\right[$ is $u - \ln s$ and its
second derivative is $-1/s < 0$, so it is maximised uniquely at $s = e^{u}$,
with value $e^{u}$; the boundary value at $s = 0$ is $0 < e^{u}$. Hence
$g^{*}(u) = e^{u}$, and by separability $f^{*}(u) = e^{u_{1}} + e^{u_{2}}$, so
$\dom f^{*} = U^{*} = \R^{2}$. Finally $x \mapsto (\ln x_{1}, \ln x_{2})$ maps
$\R^{2}_{++}$ bijectively onto $\R^{2}$, with the coordinatewise exponential as
inverse.
\end{proof}

\begin{lemma}\label{lem:C}
$C$ is nonempty and compact, $C \subseteq U$, the map $c$ is injective, and
\[
  \Cstar = \{ (t, -t^{2}) : t \in [a,b] \},
  \qquad
  \cl \Cstar = \Cstar \subseteq U^{*} .
\]
\end{lemma}

\begin{proof}
The map $c$ is continuous on the compact interval $[a,b]$, so $C$ is compact,
hence closed, and nonempty. Both coordinates of $c(t)$ are strictly positive, so
$C \subseteq \R^{2}_{++} = U$. Injectivity holds because $t \mapsto e^{t}$ is
strictly increasing. By Lemma~\ref{lem:legendre},
$\nabla f(c(t)) = (\ln e^{t}, \ln e^{-t^{2}}) = (t, -t^{2})$, so $\Cstar$ is the
stated arc; it is a continuous image of $[a,b]$, hence compact and closed, and
$U^{*} = \R^{2}$ contains it.
\end{proof}

The next lemma is the computation that makes the example tractable: the right
Bregman projection onto $C$ decouples into a one-variable problem, and the part
depending on $x$ is an additive constant.

\begin{lemma}\label{lem:reduce}
For every $x \in U$ and every $t \in [a,b]$,
\[
  \D(x, c(t)) = f(x) + h_{x}(t),
  \qquad\text{where}\qquad
  h_{x}(t) = e^{t} + e^{-t^{2}} - x_{1} t + x_{2} t^{2}.
\]
\end{lemma}

\begin{proof}
For $x, y \in U$, expanding coordinatewise gives the familiar form
\[
  \D(x,y)
  = \sum_{j=1}^{2}
    \Bigl[ x_{j} \ln x_{j} - x_{j} - (y_{j} \ln y_{j} - y_{j})
           - \ln(y_{j})(x_{j} - y_{j}) \Bigr]
  = \sum_{j=1}^{2}
    \Bigl[ x_{j} \ln \tfrac{x_{j}}{y_{j}} - x_{j} + y_{j} \Bigr].
\]
Substituting $y = c(t) = (e^{t}, e^{-t^{2}})$ and using
$\ln(x_{1}/e^{t}) = \ln x_{1} - t$ and
$\ln(x_{2}/e^{-t^{2}}) = \ln x_{2} + t^{2}$,
\[
  \D(x, c(t))
  = \bigl( x_{1} \ln x_{1} - x_{1} + x_{2} \ln x_{2} - x_{2} \bigr)
    + \bigl( e^{t} + e^{-t^{2}} - x_{1} t + x_{2} t^{2} \bigr),
\]
and the first bracket is exactly $f(x)$.
\end{proof}

\begin{lemma}\label{lem:convex}
For every $x \in U$ the function $h_{x}$ is strictly convex on $[a,b]$. Indeed
\[
  h_{x}''(t) = e^{t} + (4t^{2} - 2) e^{-t^{2}} + 2 x_{2}
  \ \ge\ e^{a} + 2 x_{2}
  \ >\ 0
  \qquad (t \in [a,b]).
\]
\end{lemma}

\begin{proof}
Differentiating $h_{x}$ twice gives
\[
  h_{x}'(t) = e^{t} - 2t e^{-t^{2}} - x_{1} + 2 x_{2} t,
  \qquad
  h_{x}''(t) = e^{t} + (4t^{2} - 2) e^{-t^{2}} + 2 x_{2} .
\]
Since $t \ge a \ge 1/\sqrt{2}$ we have $4t^{2} - 2 \ge 0$, so the middle term is
nonnegative; and $e^{t} \ge e^{a}$ while $x_{2} > 0$ because $x \in U$. The
displayed bound follows, and it is strictly positive.
\end{proof}

\begin{remark}\label{rem:sharp}
Only positivity is used. For the record, on $[1,2]$ the exact infimum of
$t \mapsto e^{t} + (4t^{2} - 2)e^{-t^{2}}$ is attained at $t = 1$ and equals
$e + 2/e = 3.45404\ldots$, so $h_{x}'' > e + 2/e$ there, uniformly in $x \in U$.
We do not need this sharper constant, and we note that establishing it requires
a genuine monotonicity argument, whereas the bound of Lemma~\ref{lem:convex} does
not.
\end{remark}

\begin{lemma}\label{lem:cheb}
$C$ is right $\D$-Chebyshev: $\Proj_{C}(x)$ is a singleton for every $x \in U$.
\end{lemma}

\begin{proof}
Fix $x \in U$. By Lemma~\ref{lem:reduce} the term $f(x)$ does not depend on $t$,
so minimising $\D(x, \cdot)$ over $C$ is the same as minimising $h_{x}$ over
$[a,b]$; and by injectivity of $c$ (Lemma~\ref{lem:C}) a unique minimising
parameter yields a unique minimising point of $C$. The function $h_{x}$ is
continuous on the compact interval $[a,b]$, so a minimiser exists. If $s \neq t$
were both minimisers then strict convexity (Lemma~\ref{lem:convex}) would give
\[
  h_{x}\!\left( \tfrac{s+t}{2} \right)
  < \tfrac{1}{2}\bigl( h_{x}(s) + h_{x}(t) \bigr)
  = \min_{[a,b]} h_{x},
\]
which is absurd. Hence the minimiser is unique.
\end{proof}

\begin{remark}\label{rem:endpoint}
The minimiser need not be interior, and no closed form for it is needed. Since
$h_{x}'$ is strictly increasing by Lemma~\ref{lem:convex}, exactly one of three
cases occurs: if $h_{x}'(a) \ge 0$ the minimiser is $t = a$; if
$h_{x}'(b) \le 0$ it is $t = b$; otherwise $h_{x}'(a) < 0 < h_{x}'(b)$ and the
minimiser is the unique zero of $h_{x}'$ in $\left]a,b\right[$. These cases are
exhaustive because $h_{x}'(a) < h_{x}'(b)$. Uniqueness in
Lemma~\ref{lem:cheb} is insensitive to which case obtains.
\end{remark}

\begin{lemma}\label{lem:nonconvex}
$\Cstar$ is not convex.
\end{lemma}

\begin{proof}
The points $(a, -a^{2})$ and $(b, -b^{2})$ lie in $\Cstar$, and their midpoint is
$m = \bigl( \tfrac{a+b}{2}, -\tfrac{a^{2}+b^{2}}{2} \bigr)$. If $m$ belonged to
$\Cstar$ its first coordinate would force the parameter $t = (a+b)/2$, and its
second coordinate would then have to equal $-\bigl( \tfrac{a+b}{2} \bigr)^{2}$.
But
\[
  \frac{a^{2}+b^{2}}{2} - \left( \frac{a+b}{2} \right)^{2}
  = \frac{(a-b)^{2}}{4} > 0 ,
\]
so $-\tfrac{a^{2}+b^{2}}{2} < -\bigl( \tfrac{a+b}{2} \bigr)^{2}$ and
$m \notin \Cstar$.
\end{proof}

\begin{proof}[Proof of Theorem~\textup{\ref{thm:main}}]
Lemma~\ref{lem:legendre} gives that $f$ is Legendre with
$\dom f = \R^{2}_{+} \neq X$, so hypothesis (a) fails. Item (i) is
Lemma~\ref{lem:C}, item (ii) is Lemma~\ref{lem:C}, item (iii) is
Lemma~\ref{lem:cheb}, and item (iv) is Lemma~\ref{lem:nonconvex}. Since (b) and
(c) hold while the conclusion of Fact~\ref{fact:BMW} fails, hypothesis (a)
cannot be deleted.
\end{proof}

\begin{remark}[hypothesis audit]\label{rem:audit}
For the reader checking that no hypothesis has been violated silently: $X$ is
finite dimensional; $f$ is Legendre; $C$ is nonempty, compact, and contained in
$U$; $\cl \Cstar \subseteq U^{*}$; and $C$ is right $\D$-Chebyshev on all of
$U$. The only failure is the intended one, $\dom f \neq X$. If one reads
Fact~\ref{fact:BMW} under an additional standing assumption of $1$-coercivity,
the example still complies: $g \ge -1$ on $[0,\infty[$, so with
$M = \max\{x_{1},x_{2}\}$ we have $f(x) \ge M \ln M - M - 1$ and
$\lVert x \rVert_{2} \le \sqrt{2}\, M$ for $x \in \dom f$, whence
$f(x)/\lVert x \rVert_{2} \to +\infty$; outside $\dom f$ the function is
$+\infty$. One should not, however, read a standing assumption of
supercoercivity for $f^{*}$ into the setting: for a closed convex $f$,
supercoercivity of $f^{*}$ is equivalent to $\dom f = X$, which is exactly the
hypothesis under test, so requiring it would make Problem~\ref{prob:2} vacuous.
\end{remark}

\begin{remark}\label{rem:C-nonconvex}
The set $C$ itself is nonconvex, though nothing above requires this;
Fact~\ref{fact:BMW} does not assume $C$ convex and concludes about $\Cstar$. In
the variable $r = e^{t}$, $C$ is the graph of
$\psi(r) = e^{-(\ln r)^{2}}$ over $[e^{a}, e^{b}]$, and with $L = \ln r$,
\[
  \psi''(r) = \frac{e^{-L^{2}}}{r^{2}} \bigl( 4L^{2} + 2L - 2 \bigr),
\]
which is strictly positive for $L \ge 1/\sqrt2$; so for $a \ge 1/\sqrt2$ the
function $\psi$ is strictly convex and $C$ is a strictly convex arc, in
particular not convex as a set.
\end{remark}

\section{A consequence for the theorem of Luo, Meng, Wen and Yao}\label{sec:luo}

Luo, Meng, Wen and Yao \cite{LMWY2019} revisit the results surveyed in
\cite{BMW2010} with the aim of removing coercivity assumptions. Their
Theorem~3.12 concerns a right $\D$-proximinal set $C$ and the statements
\begin{enumerate}[label=\textup{(\roman*)},leftmargin=2.2em]
  \item for all $x \in U$ and $y \in C$: $y \in \Proj_{C}(x)$ if and only if
        $\ip{\nabla f(y) - \nabla f(c)}{y - x} \le 0$ for every $c \in C$;
  \item $C$ is a right $\D$-sun, that is, $y \in \Proj_{C}(x)$ implies
        $y \in \Proj_{C}(\lambda x + (1-\lambda) y)$ for every $\lambda \ge 0$
        with $\lambda x + (1-\lambda) y \in U$;
  \setcounter{enumi}{3}
  \item $\nabla f(C)$ is convex.
\end{enumerate}
They prove that (i) $\Rightarrow$ (ii) with no domain hypothesis; that
(i) $\Leftrightarrow$ (ii) \emph{if $U = X$}; and that (i) $\Leftrightarrow$ (iv)
if $\nabla f(U) = U^{*}$ and $f^{*}$ is G\^ateaux differentiable and strictly
convex on $U^{*}$. Their Theorem~3.13 states that if $f$ is totally convex at
every point of $U$, if $\nabla f(U) = U^{*}$, and if $f^{*}$ is locally uniformly
totally convex at every point of $U^{*}$, then every boundedly compact right
$\D$-Chebyshev set is a right $\D$-sun.

Our example satisfies (ii) but not (i), and does not satisfy $U = X$. Both
halves can be checked directly, without invoking either of their theorems, and
we do so: the point is to exhibit a configuration their Theorem~3.12(2) would
forbid, so it is better not to rely on their machinery to produce it.

\begin{lemma}\label{lem:sun}
Let $f$ and $C$ be as in Theorem~\textup{\ref{thm:main}}. Then $C$ is a right
$\D$-sun.
\end{lemma}

\begin{proof}
Write $y = c(s) \in \Proj_{C}(x)$ for some $x \in U$ and $s \in [a,b]$, and let
$\lambda \ge 0$ satisfy $z_{\lambda} := \lambda x + (1-\lambda) y \in U$. By
Lemma~\ref{lem:reduce}, the projection of any $w \in U$ is governed by $h_{w}$,
and by Lemma~\ref{lem:convex} each $h_{w}$ is strictly convex on $[a,b]$, so
$h_{w}'$ is strictly increasing.

The essential observation is that
\[
  h_{w}'(t) = e^{t} - 2t e^{-t^{2}} - w_{1} + 2 w_{2} t
\]
is an \emph{affine} function of $w$. Moreover $h_{y}'(s) = 0$: substituting
$w = y = c(s) = (e^{s}, e^{-s^{2}})$ gives
\[
  h_{y}'(s) = e^{s} - 2s e^{-s^{2}} - e^{s} + 2 e^{-s^{2}} s = 0 .
\]
Hence, by affinity,
\[
  h_{z_{\lambda}}'(s)
  = \lambda h_{x}'(s) + (1-\lambda) h_{y}'(s)
  = \lambda\, h_{x}'(s).
\]
Now use the three cases of Remark~\ref{rem:endpoint} for $x$. If
$s \in \left]a,b\right[$ then $h_{x}'(s) = 0$, so $h_{z_{\lambda}}'(s) = 0$ and
$s$ is the unique minimiser of $h_{z_{\lambda}}$. If $s = a$ then
$h_{x}'(a) \ge 0$, so $h_{z_{\lambda}}'(a) = \lambda h_{x}'(a) \ge 0$ since
$\lambda \ge 0$, and as $h_{z_{\lambda}}'$ is increasing the unique minimiser is
$a$. If $s = b$ then $h_{x}'(b) \le 0$, so $h_{z_{\lambda}}'(b) \le 0$ and the
unique minimiser is $b$. In all three cases
$\Proj_{C}(z_{\lambda}) = \{c(s)\} = \{y\}$.
\end{proof}

\begin{lemma}\label{lem:i-fails}
Let $f$ and $C$ be as in Theorem~\textup{\ref{thm:main}}, put $s = (a+b)/2$ and
\[
  x = \bigl( e^{s} - s e^{-s^{2}},\ \tfrac{1}{2} e^{-s^{2}} \bigr) .
\]
Then $x \in U$ and $y := c(s) \in \Proj_{C}(x)$, but for every $t \in [a,b]$ with
$t \neq s$,
\[
  \ip{\nabla f(y) - \nabla f(c(t))}{y - x} = (s-t)^{2}\, \tfrac{1}{2} e^{-s^{2}} > 0 .
\]
Hence statement \textup{(i)} fails.
\end{lemma}

\begin{proof}
Both coordinates of $x$ are positive, since $e^{s} > s e^{-s^{2}}$ for
$s \ge 1/\sqrt2$, so $x \in U$. With $x_{2} = \tfrac12 e^{-s^{2}}$ we get
$2s\bigl( e^{-s^{2}} - x_{2} \bigr) = s e^{-s^{2}} = e^{s} - x_{1}$, which
rearranges to $h_{x}'(s) = e^{s} - 2s e^{-s^{2}} - x_{1} + 2 x_{2} s = 0$. As
$s \in \left]a,b\right[$ and $h_{x}$ is strictly convex, $s$ is the unique
minimiser, so $y = c(s)$ is the right projection of $x$.

By Lemma~\ref{lem:legendre}, $\nabla f(c(r)) = (r, -r^{2})$, so
\[
  \nabla f(y) - \nabla f(c(t)) = (s - t,\ t^{2} - s^{2}) = (s-t)\,(1,\ -(s+t)),
\]
while $y - x = \bigl( e^{s} - x_{1},\ e^{-s^{2}} - x_{2} \bigr)$. Therefore
\[
  \ip{\nabla f(y) - \nabla f(c(t))}{y - x}
  = (s-t) \Bigl[ (e^{s} - x_{1}) - (s+t)\bigl( e^{-s^{2}} - x_{2} \bigr) \Bigr].
\]
Substituting $e^{s} - x_{1} = 2s( e^{-s^{2}} - x_{2})$ from the stationarity
relation above, the bracket becomes
$(2s - s - t)( e^{-s^{2}} - x_{2}) = (s-t)( e^{-s^{2}} - x_{2})$, so the whole
expression equals $(s-t)^{2}( e^{-s^{2}} - x_{2}) = (s-t)^{2} \cdot
\tfrac12 e^{-s^{2}}$, which is strictly positive for $t \neq s$. Since $a < b$
such a $t$ exists in $[a,b]$. Statement (i) requires this quantity to be
$\le 0$ for every $c \in C$, so (i) fails.
\end{proof}

\begin{corollary}\label{cor:luo}
Let $f$ and $C$ be as in Theorem~\textup{\ref{thm:main}}. Then \textup{(ii)}
holds, \textup{(i)} and \textup{(iv)} fail, and $U \neq X$. Consequently the
hypothesis $U = X$ in \textup{\cite[Theorem 3.12(2)]{LMWY2019}} cannot be
deleted without replacement: the implication \textup{(ii)} $\Rightarrow$
\textup{(i)} is false without it.
\end{corollary}

\begin{proof}
Statement (ii) is Lemma~\ref{lem:sun}, statement (i) fails by
Lemma~\ref{lem:i-fails}, and (iv) fails by Lemma~\ref{lem:nonconvex}. Finally
$U = \R^{2}_{++} \neq \R^{2} = X$.
\end{proof}

We state the conclusion as non-omissibility rather than necessity. The example
shows that $U = X$ cannot simply be struck from Theorem~3.12(2); it does not
show that $U = X$ is the weakest hypothesis that would serve there, and some
weaker ray-completeness condition may well suffice.

\begin{remark}[where the obstruction lies]\label{rem:obstruction}
The example locates the failure precisely. The proof of
(ii) $\Rightarrow$ (i) in \cite[Theorem 3.12(2)]{LMWY2019} applies the sun
property along $z_{\lambda}$ for arbitrarily large $\lambda$, divides by
$\lambda$, and lets $\lambda \to \infty$. For the witness of
Lemma~\ref{lem:i-fails} that limit is unavailable: with
$\eta = e^{-s^{2}}$ we have $(z_{\lambda})_{2} = \lambda \tfrac{\eta}{2} +
(1-\lambda)\eta = \eta\bigl( 1 - \tfrac{\lambda}{2} \bigr)$, so
$z_{\lambda} \in U$ exactly for $0 \le \lambda < 2$. The ray leaves $U$ after a
bounded distance, and it is this bounded reach, not the sun property itself,
that $U = X$ was supplying.
\end{remark}

\begin{remark}
No result of \cite{LMWY2019} is contradicted. Their Theorem~3.12(1) gives
(i) $\Rightarrow$ (ii) unconditionally, which is consistent with (ii) holding and
(i) failing. Their Theorem~3.12(3) gives (i) $\Leftrightarrow$ (iv) under
hypotheses our example does satisfy, and indeed (i) and (iv) fail together here,
as Lemmas~\ref{lem:nonconvex} and~\ref{lem:i-fails} confirm independently. Their
Theorem~3.13 also applies to our example, by a routine verification that
negative entropy is totally convex on $U$ and that $f^{*}$ is locally uniformly
totally convex on $U^{*}$; it yields Lemma~\ref{lem:sun} again, and we gave the
direct proof instead because it is shorter and does not depend on that
verification. Their Corollary~3.14, which does conclude convexity of
$\nabla f(C)$ from bounded compactness and right Chebyshevness, assumes total
convexity at every point of $X$ and $\nabla f(X) = U^{*}$; both presuppose
$U = X$ and therefore do not apply here.
\end{remark}

\begin{remark}
It is instructive that the same example does duty twice. Both
\cite[Fact 3.2]{BMW2010} and \cite[Theorem 3.12(2)]{LMWY2019} reach convexity of
$\nabla f(C)$ from a Chebyshev-type hypothesis, and both require full domain at
the corresponding step, one by assuming $\dom f = X$ outright and the other by
requiring $U = X$. Theorem~\ref{thm:main} and Corollary~\ref{cor:luo} together
show the requirement is not an artefact of either proof technique.
\end{remark}

\section{The remaining hypothesis}\label{sec:b}

For completeness we record the status of hypothesis (b) of
Fact~\ref{fact:BMW}, which is not ours to claim.

\begin{proposition}[essentially {\cite[Theorems 3.12 and 3.13]{LMWY2019}}]
\label{prop:b}
Let $f$ be Legendre on $X = \R^{n}$ with $\dom f = X$, and let $C \subseteq X$
be nonempty and closed. If $C$ is right $\D$-Chebyshev, then $\Cstar$ is convex.
In particular hypothesis \textup{(b)} of Fact~\textup{\ref{fact:BMW}} is
redundant.
\end{proposition}

\begin{proof}
Since $\dom f = X$ we have $U = X$, and $C$, being closed in $\R^{n}$, is
boundedly compact. By \cite[Theorem 26.5]{Rock1970}, $\nabla f$ is a bijection
of $U = X$ onto $U^{*}$ with inverse $\nabla f^{*}$, and $f^{*}$ is Legendre; in
particular $\nabla f(U) = U^{*}$, and $f^{*}$ is differentiable and strictly
convex on $U^{*}$.

The total-convexity hypotheses of \cite[Theorem 3.13]{LMWY2019} are automatic
here. For $y \in X$ and $t > 0$ the set $\{z \in \dom f = X :
\lVert z - y \rVert = t\}$ is compact, $\D(\cdot, y)$ is continuous on it, and
$\D(z,y) > 0$ for $z \neq y$ by strict convexity, so
$\nu_{f}(y,t) > 0$. For $f^{*}$, fix $u_{0} \in U^{*}$ and $t > 0$ and choose
$\rho > 0$ with $\overline{B}(u_{0},\rho) \subseteq U^{*}$; the \emph{joint}
set $K = \{(u,v) : \lVert u - u_{0}\rVert \le \rho,\ \lVert v - u \rVert = t\}$
is compact and contains no pair with $v = u$, so $m := \min_{K} D_{f^{*}} > 0$
and $\nu_{f^{*}}(u,t) \ge m$ for every $u \in \overline{B}(u_{0},\rho)$, whence
$\nu^{\mathrm{loc}}_{f^{*}}(u_{0},t) \ge m > 0$. (A bound for each fixed $u$
would not suffice, since such bounds could degenerate as $u \to u_{0}$.)

So \cite[Theorem 3.13]{LMWY2019} gives that $C$ is a right $\D$-sun, which is
(ii); \cite[Theorem 3.12(2)]{LMWY2019}, applicable since $U = X$, yields (i);
and \cite[Theorem 3.12(3)]{LMWY2019} makes (i) equivalent to convexity of
$\nabla f(C)$.
\end{proof}

Together with Theorem~\ref{thm:main}, this completes the account of
Fact~\ref{fact:BMW}: hypothesis (a) is essential, and hypothesis (b) is not
needed.

\section{What remains open}\label{sec:open}

Theorem~\ref{thm:main} answers Problem~\ref{prob:2} as posed, but leaves the
more interesting half of the question untouched.

\begin{enumerate}[leftmargin=2.2em]
\item \emph{What replaces full domain?} Our example shows $\dom f = X$ cannot
simply be deleted. It does not identify the weakest hypothesis under which the
conclusion survives. A natural guess is a condition on how $\Cstar$ approaches
$\partial U^{*}$, since by the duality of
\cite[Lemma 3.2]{LMWY2019} the right projection onto $C$ for $f$ is the left
projection onto $\Cstar$ for $f^{*}$, and the obstruction in our example lives
entirely on the dual side.

\item \emph{Problem 4 of \cite{BMW2010}} asks for a characterisation of the
right $\D$-Chebyshev sets for the negative entropy. A solution would subsume
Theorem~\ref{thm:main}, which exhibits one such set and computes its gradient
image. We regard this as the natural next question, and note that our example
is a member of a one-parameter family (any $1/\sqrt2 \le a < b$), which suggests
such sets are not scarce.

\item \emph{Dimension and geometry.} Our example is planar and our curve is an
arc. We do not know whether the failure can be made worse, for instance with
$\Cstar$ having nonempty interior in its convex hull, nor what the analogous
picture is in $\R^{n}$ for $n \ge 3$.
\end{enumerate}

\section*{Verification}

The claims of Theorem~\ref{thm:main} are elementary and were checked by hand,
but they were also checked numerically before and after the proofs were written.
A harness verifies, in directed-rounding interval arithmetic and to $40$ digits:
that $\D$ agrees with its definition and with the generalised
Kullback--Leibler form; that the reduction of Lemma~\ref{lem:reduce} holds with
the constant equal to $f(x)$; that $h_{x}''$ is bounded below uniformly on
$[1,2]$; that the minimiser of $h_{x}$ is unique for a spread of adversarial
$x \in U$ including near-boundary and far-field points; and the nonconvexity
witness of Lemma~\ref{lem:nonconvex}. A second harness tests
Lemma~\ref{lem:sun} directly against the definition of a sun, over $1{,}038$
pairs $(x, \lambda)$ with $\lambda$ ranging beyond $x$ to the boundary of $U$
and including points constructed to project to the interior of the arc; the
projection parameter is invariant to $4 \times 10^{-40}$. It also confirms the
witness of Lemma~\ref{lem:i-fails}, both the closed form
$(s-t)^{2}\bigl(e^{-s^{2}} - x_{2}\bigr)$ for the inner product and its
positivity. These checks are finite and do not certify the universal
quantifiers, which are carried by the proofs above.

\begin{thebibliography}{9}

\bibitem{BBC2001}
H.~H. Bauschke, J.~M. Borwein, and P.~L. Combettes,
\emph{Essential smoothness, essential strict convexity, and Legendre functions
in Banach spaces},
Commun. Contemp. Math. \textbf{3} (2001), 615--647.

\bibitem{BWYY2009}
H.~H. Bauschke, X.~Wang, J.~Ye, and X.~Yuan,
\emph{Bregman distances and Chebyshev sets},
J. Approx. Theory \textbf{159} (2009), no.~1, 3--25.

\bibitem{BMW2010}
H.~H. Bauschke, M.~S. Macklem, and X.~Wang,
\emph{Chebyshev sets, Klee sets, and Chebyshev centers with respect to Bregman
distances: recent results and open problems},
\href{https://arxiv.org/abs/1003.3127}{arXiv:1003.3127} (2010).

\bibitem{LOC2020}
E.~Laude, P.~Ochs, and D.~Cremers,
\emph{Bregman proximal mappings and Bregman--Moreau envelopes under relative
prox-regularity},
J. Optim. Theory Appl. \textbf{184} (2020), 724--761.

\bibitem{LMWY2019}
X.-F. Luo, K.~Meng, C.-F. Wen, and J.-C. Yao,
\emph{Bregman distances without coercive condition: suns, Chebyshev sets and
Klee sets},
Optimization \textbf{68} (2019), no.~8, 1599--1624.

\bibitem{TW2025}
A.~Themelis and Z.~Wang,
\emph{On the natural domain of Bregman operators},
\href{https://arxiv.org/abs/2506.00465}{arXiv:2506.00465} (2025).

\bibitem{Rock1970}
R.~T. Rockafellar,
\emph{Convex Analysis},
Princeton University Press, 1970.

\end{thebibliography}

\end{document}
